% MathorCup 2026 大数据竞赛 — 论文模板
% 编译：make paper   （内部走 xelatex，中文用 ctex）
% 结构：标题/摘要 → 正文 → 参考文献 → 附录（代码）
\documentclass[12pt, a4paper]{ctexart}

\usepackage{geometry}
\geometry{left=2.5cm, right=2.5cm, top=2.5cm, bottom=2.5cm}
\usepackage{amsmath, amssymb}
\usepackage{graphicx}
\usepackage{booktabs}      % 三线表 \toprule \midrule \bottomrule
\usepackage{multirow}
\usepackage{siunitx}       % 数字对齐
\usepackage{caption}
\usepackage{float}
\usepackage{hyperref}
\usepackage[numbers]{natbib}   % 参考文献，配合 refs.bib（\citep）

\graphicspath{{figures/}}
\captionsetup{font=small, labelfont=bf}

% ---- 标题信息：开赛拿到赛题后改这里 ----
\title{\textbf{题目：〔赛题名称〕}\\[6pt]
       \large ——〔副标题：一句话概括你的方法〕}
\author{}
\date{}

\begin{document}

% ============ 标题页 ============
\maketitle
\thispagestyle{empty}

\vspace{-1cm}
\begin{center}
\begin{tabular}{ll}
\toprule
参赛队号 & 〔报名后官网生成，填这里〕 \\
参赛组别 & 〔研究生组 / 本科组 / 专科组〕 \\
所选赛题 & 〔A 题 / B 题〕 \\
\bottomrule
\end{tabular}
\end{center}

\vspace{0.5cm}

% ============ 摘要 ============
% 评审明确要求：摘要要"简明扼要地概述解决问题的核心思想以及效果"。
% 写法建议：问题一句 → 方法一段 → 结果用数字 → 结论一句。
\begin{center}
\textbf{\large 摘\quad 要}
\end{center}

\noindent
〔一句话点出赛题要解决的核心问题。〕
针对〔问题一〕，本文〔核心方法，含关键创新〕，得到〔主结果，写具体数字，
如：准确率 0.923、较基线提升 4.1 个百分点〕。
针对〔问题二〕，〔方法〕，〔结果〕。
本文的主要创新点在于：〔1. …；2. …；3. …〕。
最后，〔结论与推广价值，一句话〕。

\vspace{0.3cm}
\noindent\textbf{关键词：}〔关键词1；关键词2；关键词3；关键词4〕

\newpage

% ============ 正文 ============
\section{问题重述与分析}
\subsection{问题背景}
\subsection{问题提出}
\subsection{问题分析}   % 三问逐条拆解，画出思路流程图

\section{数据预处理与探索性分析}
\subsection{数据概况}
\subsection{数据清洗}   % 缺失、异常、重复；说明每条处理规则的理由
\subsection{探索性分析} % 放图：分布、相关性、类别不平衡
% 示例分节文件。main.tex 里用 % 示例分节文件。main.tex 里用 % 示例分节文件。main.tex 里用 \input{sections/eda} 引入。
% 好处：多人各写一节，减少冲突；单节编译调试快。

数据基本情况见表~\ref{tab:data-overview}。

\begin{table}[H]
\centering
\caption{数据概况}
\label{tab:data-overview}
\begin{tabular}{lrrr}
\toprule
数据集 & 样本数 & 特征数 & 缺失率 \\
\midrule
训练集 & —— & —— & —— \\
测试集 & —— & —— & —— \\
\bottomrule
\end{tabular}
\end{table}

由图~\ref{fig:example} 可见……

\begin{figure}[H]
\centering
% \includegraphics[width=0.8\textwidth]{example.png}
\fbox{\parbox[c][4cm][c]{0.8\textwidth}{\centering （此处放图）}}
\caption{〔图题：描述这张图说明了什么，而不是"某某分布图"〕}
\label{fig:example}
\end{figure}
 引入。
% 好处：多人各写一节，减少冲突；单节编译调试快。

数据基本情况见表~\ref{tab:data-overview}。

\begin{table}[H]
\centering
\caption{数据概况}
\label{tab:data-overview}
\begin{tabular}{lrrr}
\toprule
数据集 & 样本数 & 特征数 & 缺失率 \\
\midrule
训练集 & —— & —— & —— \\
测试集 & —— & —— & —— \\
\bottomrule
\end{tabular}
\end{table}

由图~\ref{fig:example} 可见……

\begin{figure}[H]
\centering
% \includegraphics[width=0.8\textwidth]{example.png}
\fbox{\parbox[c][4cm][c]{0.8\textwidth}{\centering （此处放图）}}
\caption{〔图题：描述这张图说明了什么，而不是"某某分布图"〕}
\label{fig:example}
\end{figure}
 引入。
% 好处：多人各写一节，减少冲突；单节编译调试快。

数据基本情况见表~\ref{tab:data-overview}。

\begin{table}[H]
\centering
\caption{数据概况}
\label{tab:data-overview}
\begin{tabular}{lrrr}
\toprule
数据集 & 样本数 & 特征数 & 缺失率 \\
\midrule
训练集 & —— & —— & —— \\
测试集 & —— & —— & —— \\
\bottomrule
\end{tabular}
\end{table}

由图~\ref{fig:example} 可见……

\begin{figure}[H]
\centering
% \includegraphics[width=0.8\textwidth]{example.png}
\fbox{\parbox[c][4cm][c]{0.8\textwidth}{\centering （此处放图）}}
\caption{〔图题：描述这张图说明了什么，而不是"某某分布图"〕}
\label{fig:example}
\end{figure}
   % 大段内容拆到 sections/ 下，模板见该文件

\section{模型建立与求解}
\subsection{特征工程}
\subsection{模型选择与原理}
\subsection{训练与调参}
\subsection{结果与分析}

\section{模型评价与推广}
\subsection{模型优点}
\subsection{模型缺点}
\subsection{改进方向}

\section{结论}

% ============ 参考文献 ============
\newpage
\bibliographystyle{plainnat}
\bibliography{refs}

% ============ 附录 ============
\newpage
\appendix
\section{程序代码}
% 注意：即使代码放进了这里，官方要求仍须单独以支撑材料提交一份源程序。
\subsection{数据预处理}
\subsection{模型训练}

\end{document}
